\documentclass[12pt]{article}
\usepackage{graphicx}
\usepackage[margin=1in]{geometry}
\usepackage{caption}
\usepackage{xcolor}
\usepackage{float}
\usepackage{amsmath}
\usepackage{array}
\usepackage{multicol}
\usepackage{listings}
\usepackage{hyperref}
\usepackage{url}
\usepackage{fancyvrb}

\setlength{\columnsep}{1cm}

\definecolor{codegray}{gray}{0.95}
\lstset{
    backgroundcolor=\color{codegray},
    basicstyle=\ttfamily\small,
    columns=fullflexible,
    frame=single,
    breaklines=true,
    postbreak=\mbox{\textcolor{red}{$\hookrightarrow$}\space},
    showstringspaces=false
}

\begin{document}

% Title Block
\begin{figure}[H]
    \begin{minipage}{0.45\textwidth}
        \includegraphics[width=\textwidth]{WhatsApp Image 2025-07-08 at 3.59.18 PM.jpeg} % <-- Replace with actual image if available
    \end{minipage}
\hfill
    \begin{minipage}{0.45\textwidth}
        \textbf{Name: Madhu Latha G} \\
        \textbf{ID: cometfwc029} \\
        \textbf{Date: 9th July 2025}
    \end{minipage}
\end{figure}

\begin{center}
    {\LARGE \textbf{\textcolor{cyan}{GATE Question Paper 2010, EE Question Number 54}}}
\end{center}
\vspace{1em}

\begin{multicols}{2}

% Abstract
\noindent\textbf{Abstract} \\[0.5em]
\includegraphics[width=0.99\linewidth]{Screenshot from 2025-07-09 16-15-09.png} \\[0.5em]
\textit{(GATE 2010, Question No. 54 – Natural Response of LC Circuit: Simulated on Raspberry Pi Pico)} \\

\vspace{1em}
\noindent\textbf{1. Components Used}
\begin{table}[H]
\small
\centering
\begin{tabular}{|p{4.2cm}|c|}
\hline
\textbf{Component} & \textbf{Qty} \\
\hline
Raspberry Pi Pico (RP2040) & 1 \\
Breadboard & 1 \\
LED & 1 \\
220$\Omega$ Resistor & 1 \\
Jumper Wires (M-M) & 6 \\
USB Cable (Micro-B) & 1 \\
Computer with Thonny IDE & 1 \\
\hline
\end{tabular}
\caption*{Table 1: List of components used}
\end{table}

\vspace{1em}
\noindent\textbf{2. Setup and Connections}
\begin{enumerate}
    \item Connect GP15 (PWM pin) to LED through 220$\Omega$ resistor.
    \item Connect cathode of LED to GND.
    \item Connect Raspberry Pi Pico to PC using USB cable.
    \item Flash MicroPython firmware and open Thonny IDE.
\end{enumerate}

\vspace{1em}
\noindent\textbf{3. Simulation Strategy}
\begin{itemize}
    \item Generate a cosine waveform: $i(t) = 10 \cos(5000t)$.
    \item Use MicroPython to simulate waveform using PWM.
    \item LED brightness varies according to the waveform.
\end{itemize}

\vspace{1em}

\vspace{1em}
\noindent\textbf{5. Working Principle}
\begin{itemize}
    \item The waveform simulates current in an LC circuit.
    \item The LED acts as a visual indicator of waveform magnitude.
    \item Cosine waveform is generated using a lookup table.
    \item 100 samples/cycle approximates the waveform smoothly.
\end{itemize}

\vspace{1em}
\noindent\textbf{6. Observations and Output}
\begin{itemize}
    \item LED brightness oscillates sinusoidally.
    \item Represents undamped oscillations in LC circuit.
    \item Waveform can be viewed on oscilloscope from GP15.
\end{itemize}

\vspace{1em}
\noindent\textbf{7. Result}
The waveform $i(t) = 10 \cos(5000t)$ is successfully emulated using PWM on Raspberry Pi Pico. It matches the expected LC circuit response.

\vspace{1em}
\noindent\textbf{8. Source Code Link} \\
\url{https://github.com/madhuu34/Madhulathafwc/blob/main/Hardware/AVR_GCC/hardware_code.py}

\end{document}

